\chapter{Conclusion}
\label{ch:conclusion}

\section{Answers to the Research Questions}

This thesis asked two questions: how a frequent batch auction engine is
designed, and what changes when the same real order flow is cleared by
one instead of by a continuous book. The two paragraphs below answer each
in turn.

\paragraph{RQ1.} A frequent batch auction engine is fully specified by five
design blocks: the batch interval, the disclosure regime during
accumulation, the order flow segregation architecture, the clearing and
rationing rule, and settlement and residual handling. These are not
optional extras layered onto the mechanism but decisions a continuous
engine also makes implicitly, which become visible and contestable only
once serial processing is abandoned. Three points carry the most weight.
The only unambiguous structural gain is discrete accumulation combined with
uniform pricing; every other choice trades one distortion for another, so
the question is which distortion the venue prefers. The blocks interact and
cannot be optimised independently. And a venue can implement frequent batch
auctions exactly as proposed and still have a speed race, because
price-time rationing inside the batch reintroduces one through the
allocation rule --- the sharpest practical implication of the analysis.
Two questions with no settled answer were developed as contributions: the
maker--taker distinction dissolves under simultaneous uniform-price
clearing, with persistence across batches the strongest of four candidate
redefinitions; and the price improvement uniform clearing produces is not a
fixed property of the mechanism but depends on the clearing and rationing
rules chosen.

\paragraph{RQ2.} The paired comparison of the two engines on identical
order flow covers the one configuration actually run: SOL/USD,
$\tau=1$\,second, the full month of December 2025, without the interval
sweep or rationing comparison the fuller design calls for. One finding
does most of the work in this answer: four metrics that measure the cost
of trading from different angles --- effective spread, realized spread, price impact and
Kyle's lambda --- are all close to zero for the FBA and all clearly costly
for the CDA. In the CDA, a resting order is predictably picked off: the
price moves against it after every trade, which is the textbook signature
of adverse selection. Anchoring the FBA's clearing price to its own
previous value removes that exposure, because a batch leaves no resting
order sitting there to be picked off in the first place. (The FBA has no
directly observable aggressor side, so three of the four rest on a shared
tick-rule proxy, not a certainty; Kyle's lambda needs no such proxy and
agrees with them, which is what makes the result credible despite that.)

This gain is not free, and it does not come from trading less or at a
different price --- both engines clear statistically the same volume, at
prices that track each other closely. What it costs instead is depth and
time. Depth within a fixed distance of the reference price, the one depth
measure genuinely comparable across the two engines, is slightly
\emph{lower} for the FBA at every threshold, and a given amount of trading
moves the FBA's price marginally more, not less (Amihud illiquidity) ---
the opposite of what a liquidity-pooling argument for batching would
predict, if a more modest effect than the spread result above. Execution
also simply takes longer under the FBA. Underlying both is a market that,
at this interval, is usually not the well-populated auction the mechanism
assumes: \FbaEmptyBatchShare\% of one-second batches clear no trade at all,
and even among those that do, most of the longer side is typically left
unmatched at the clearing price. At
$\tau=1$\,second on this instrument, the FBA behaves closer to its own
degenerate, near-continuous limit than to a well-filled multilateral
auction --- which is the frame every other number above has to be read
through.

\section{Limitations}
\label{sec:limitations}

The central limitation is intrinsic to the design. A replay holds order flow
constant, which is what makes the comparison clean, and by the same token
prevents participants from responding to the mechanism: the orders replayed
were submitted under a continuous engine, so their sizes were not chosen
against the rationing rule implemented (price-time-within-the-batch, a
different rule from the pro rata and size-priority alternatives surveyed
in \Cref{sec:rationing}), their arrival times
were not chosen against a batch boundary, and their limit prices were not
chosen under uncertainty about a clearing price. Every equilibrium
prediction of the FBA literature --- narrower quoted spreads, size
inflation, boundary concentration, bid shading --- is therefore out of
reach, and the simulation delivers the mechanical baseline against which
such responses would be measured. Two further limitations bound what future
work can build on: the engines are per-pair, so the multi-asset clearing
problem is part of the design analysis but not of the measurement; and
several metrics require adaptation to a batch setting, most notably the
sign of the effective and realized spread, so an extension inherits that
adaptation rather than a clean continuous-market baseline
(\Cref{sec:metricreliability} works through this case by case).

The sample itself bounds the claims further. The data come from one
venue, one instrument and one continuous month; a crypto perpetuals
exchange has a participant population, fee structure and latency profile
that differ from an equity market, so magnitudes should not be carried
over elsewhere without re-estimation, and the month gives day-by-day
stability within the sample, not a cross-instrument or cross-regime one.
And because the replay does not enforce time-in-force, it trades about
\VerCdaOverVenue\ times the venue's own recorded volume
(\Cref{sec:testfidelity}). Both engines share this gap identically, so
the paired comparisons in \Cref{ch:results} are unaffected by it. But the
absolute volume, trade-count and latency levels reported there describe
this replay, not Hyperliquid's live market.
